\section{The arithmetic picture and certified finite checks}\label{sec:diagram}
\subsection{A phase diagram with phase conditions}
Put $x=2\beta/L$. Denote the upper Abel-averaged fastest exponent at phase $\theta$ by $\alpha_u^+(\theta)$ and its supremum over phases by $\alpha_{\rm worst}$. Paper~C also defines an exponent uniform over launch sites, $\alpha_u^{\rm unif}$, which bounds that supremum. These quantities agree in the exact-law regime specified below.

The typical fastest exponent and the two packing dimensions of every nonzero spectral measure are $\max\{0,1-2/x\}$ at almost every phase, by Paper~A~\cite[Theorems~A and C]{PaperA}; the value is zero at $x=0$. The packing entries concern spectral measures, not the packing dimension of the spectral set. At $\beta=\infty$ these values are one at every phase. Paper~A gives the finite-positive-$\beta$ law used here. Its $\beta=\infty$ transport and packing endpoint is due to Jitomirskaya--Zhang~\cite[Corollary~3]{JZ2022}, and its $\beta=0$ fastest endpoint to Damanik--Tcheremchantsev~\cite{DT2007}. Paper~A sharpens earlier packing estimates of Jitomirskaya--Liu--Tcheremchantsev~\cite{JLT2025}. The typical transport framework for general supercritical type~I operators is developed in Paper~B~\cite{PaperB}.

Table~\ref{tab:phase} keeps the phase quantifiers visible. ``AL'' means Anderson localization and ``SC'' purely singular continuous spectrum. Completely resonant phases satisfy $2\theta\in\alpha\Z+\Z$. The phase $1/4$ is not completely resonant for irrational $\alpha$.

\begin{table}[ht]
\centering\footnotesize
\setlength{\tabcolsep}{4pt}
\begin{tabular}{@{}>{\raggedright\arraybackslash}p{.08\linewidth}>{\raggedright\arraybackslash}p{.23\linewidth}>{\raggedright\arraybackslash}p{.26\linewidth}>{\raggedright\arraybackslash}p{.23\linewidth}>{\raggedright\arraybackslash}p{.12\linewidth}@{}}
\toprule
$x$ & Spectral type at a.e. phase & Specified exceptional phases & $\alpha_{\rm worst}$ & $\rho$\\
\midrule
$0$ & AL & AL at completely resonant phases; SC at a dense set of other phases & $0$ & $1$ if $\alpha$ is Diophantine\\[3pt]
$(0,1)$ & AL & AL at completely resonant phases & $[0,x/(x+2)]$ & $[0,1-x]$; $\rho\ge1/C$ if $x\le2/C$\\[3pt]
$1$ & AL & No equality-line classification from the cited strict criterion & $[0,1/3]$ & Open\\[3pt]
$(1,2)$ & AL & SC at $1/4$ if $\beta_e=\beta$ & $[1-1/x,x/(x+2)]$ if $\beta_e=\beta$ & $0$\\[3pt]
$2$ & Not classified by the strict arithmetic transition & SC at $1/4$ if $\beta_e=\beta$ & $1/2$ if $\beta_e=\beta$ & $0$\\[3pt]
$(2,\infty)$ & SC at every phase & Same all-phase statement & $1-1/x$ if $\beta_e=\beta$ & $0$\\[3pt]
$\infty$ & SC at every phase & Same all-phase statement & $1$ & $0$\\
\bottomrule
\end{tabular}
\caption{Established statements and explicit open entries. Every bracket in the transport column is a pair of proved bounds. The gap column uses Theorems~\ref{thm:below}, \ref{thm:polynomial} and~\ref{thm:measure}, and Corollary~\ref{cor:bracket}. No exact gap law is asserted at $x=1$.}\label{tab:phase}
\end{table}

The a.e.-phase AL and all-phase SC statements use Avila--You--Zhou~\cite[Theorem~1.1]{AYZ2017} for $0<\beta<\infty$, with strict inequalities at $L=\beta$. The endpoint spectral-type assertions at $x=0$ and $x=\infty$ use the classical endpoint results recalled in Avila--You--Zhou (Remark~1.2 of arXiv:1512.03124). Liu~\cite[Theorem~1.1]{Liu2022} proves completely resonant AL for $L>2\beta$. At $1/4$, even $q_n$ imply a phase resonance at $q_n/2$ with error $|\delta_n|/2$; its exponent is at least $2\beta_e$. The absence-of-point-spectrum part of Jitomirskaya--Liu~\cite[Theorem~4.2 and the remark after Theorem~1.1]{JLphase} therefore gives SC when $L<2\beta_e$. Its localization half has additional frequency hypotheses; we do not use that half for this entry.

The transport bounds are from Paper~C, Theorem~7.5 and Corollaries~5.6 and~7.6. The phase-$1/4$ lower bound uses even denominators and $\beta_e$, not even indices. Without $\beta_e=\beta$, retain that lower bound with $\beta_e$ and the typical all-phase lower bound separately. The exact worst-phase formula is then not supplied by these imports. A conditional parity-free mechanism in Paper~C requires its own wide-band hypothesis; it is not a consequence of the gap theorem here.

At a finite scale,
\begin{equation}\label{eq:finitecross}
 \frac{\log\kappa_n}{Lq_n}
 =1-\frac{2b_n}{L}+\frac{2\log(2\pi q_n)}{Lq_n}.
\end{equation}
Thus $x=1$ does not force $\kappa_n=O(1)$. Fixed margins on either side of one give exponential smallness or largeness along the corresponding scales. The logarithmic correction matters at the equality line.

\subsection{The known transport geometry}
For even $q$, put $\sigma=(-1)^{q/2}$ and let $0<\eta\le1/4$. In one period of the local phase drift, the launch set
\[
 \left|\cos2\pi(q\theta+j\delta)-\sigma\right|
                     \le\eta\Lambda^{-1}/16
\]
has cardinality $q\sqrt{\eta/(2\kappa)}(1+O(\Lambda^{-1}))+O(1)$. This is the quadratic turning-point width divided by the mesh $|\delta|$. Its logarithm has rate $(b-L/2)q$ only when the main count diverges. Continuous window length is not an exact test for occupancy by one lattice cell.

Paper~C's nearby-phase burst theorem applies with its scale condition $b\ge h/2+(h-L)$ and time interval $e^{(h/2+(h-L))q}\le T\le|\delta|^{-1}$. It gives a polynomially small fastest tail. At the resonant time its radius has exponent $b-h/2$, approaching the preceding geometric exponent as $h\downarrow L$.

The confinement estimate of Paper~C, Proposition~6.4, is a uniform square-root upper bound on a cosine sublevel arc, plus its collars. At a generic phase a particular arc can be narrower. In Paper~A's notation $J_n=3+C_hq^{5/2}e^{hq}|\delta|$ is greater than three; it is $3+o(1)$ on $b\ge L+\gamma$ if $h-L<\gamma$. Alternatively $\|S^q+S^{-q}+R_q\|\le3$ eventually there. These bounds match the turning-point exponential radius. Below $b=L$, the coarse norm bound permits a larger confinement upper bound; it does not prove that this enlargement is attained by the dynamics.

\subsection{What the numerical checks certify}\label{sec:numerics}
All numbers in this subsection are inherited from the independent verification reports D77 and D80. They are interval certificates or exact finite-matrix counts. We do not rerun those campaigns here. No number is used in a proof.

For the Mathieu gaps $G_D(\kappa)$ of Appendix~\ref{app:mathieu}, D77 encloses a crossing $G_1=2G_2$ in
\[
 \kappa_*\in[1.4260784738\mathbin{\pm}6.60\cdot10^{-11}].
\]
This certifies a crossing, not global uniqueness. Ten independent checks of the rational phase-union prediction give the intervals in Table~\ref{tab:collapse}. Here the parent is $p/q$, $A$ is the following continued-fraction coefficient in the verified family, $w$ is the measured rational width, and $P'$ is the parent derivative at the measured location. This tests widths conditional on locations.

\begin{table}[ht]
\centering\small
\begin{tabular}{@{}rrrrl@{}}
\toprule
Parent & $\lambda$ & $A$ & $|m|$ & Enclosure of $\log(w|P'|/G_D)$\\
\midrule
$2/3$ &5&699&3&$[0.0001979253,0.0001979254]$\\
$3/5$ &5&2000&10&$[0.0000162096,0.0000162097]$\\
$3/5$ &2&600&15&$[-0.0005267205,-0.0005267204]$\\
$8/13$&5&385&78&$[0.0039697174,0.0039697175]$\\
$8/13$&2&77&52&$[0.0172042988,0.0172042989]$\\
$5/8$ &2&200&16&$[0.0003463309,0.0003463310]$\\
$3/5$ &5&60&15&$[0.0091106244,0.0091106245]$\\
$2/3$ &5&69&3&$[0.0034522615,0.0034522616]$\\
$8/13$&2&385&13&$[0.0005374229,0.0005374230]$\\
$2/3$ &5&149&6&$[0.0008759784,0.0008759785]$\\
\bottomrule
\end{tabular}
\caption{Certified finite rational Mathieu collapse from D77, Section~5. There are ten cases and eight distinct integer labels. These data do not prove an irrational gap law.}\label{tab:collapse}
\end{table}

D80 tests the translated construction itself on boxes containing several wells. All three cases have parent $5/13$. Case G uses $\lambda=2$, $\theta=0.0377$, $\hbar=0.1$ and three wells. H uses the same coupling and phase, $\hbar=0.005$ and two wells. OOS was registered before its independent measurement: $\lambda=3$, $\theta=0.0193$, $\hbar=0.1$ and three wells. The arithmetic uses $\delta=\hbar/(2\pi q\lambda^{q/2})$.

\begin{table}[ht]
\centering\small
\begin{tabular}{@{}llcl@{}}
\toprule
Case & Inclusive box & Counts for $D=1,\ldots,6$ & Maximum level error\\
\midrule
G & $[731,222520]$ & $3,6,9,12,15,18$ & $[0,0.000282844609]$\\
H & $[14638,2971820]$ & $2,4,6,8,10,12$ & $[0,0.000314943324]$\\
OOS & $[256912,3351000]$ & $3,6,9,12,15,18$ & $[0,0.000013630460]$\\
\bottomrule
\end{tabular}
\caption{D80 finite-box certificates. Errors are in units of $2\hbar/P_*'$ after the common dipole shift. Each row includes three interior fractions of each counting window and midpoint counts after moving both ends by $-1000$ or $+1000$ sites.}\label{tab:globalchecks}
\end{table}

The ground-cluster projector assigns each translated anchor class the per-well share
\[
\begin{split}
 \text{G: }&[0.9558209196836182962179037231977
                 \mathbin{\pm}1.93\cdot10^{-32}],\\
 \text{OOS: }&[0.98136338757751036956774744137
                 \mathbin{\pm}3.53\cdot10^{-30}].
\end{split}
\]
The tables and the display above report the independent certificates, including the corrected cluster-projector weights. The original single-eigenvector weight prediction, residual-growth prediction and factor-30 switch prediction for the ratio of a gap to the derivative proxy TP1 failed; TP1 is not a measured transport exponent. At $q=21$ and $\lambda=5$, the literal TP1 interval also failed, and the full original gap-switch calculation remains uncertified. These finite checks do not establish asymptotic complementarity. The accompanying Numerical Record, \S\S1--3, preserves the registrations, their certified comparison intervals, the corrected projector interpretation and the still uncomputed full-period test.

\subsection{New finite certificates below the threshold}\label{sec:below-numerics}
The three new independent checking campaigns certify finite rational widths, finite rounded-span Green columns, and finite rounded-span complementary spectra. These are different experiments. They do not certify exact infinite-support Wannier functions or an irrational asymptotic theorem. The accompanying Numerical Record, \S\S4--6, gives the fuller intervals and preserves failed registrations. No new computation was used to produce the intervals printed here.

The rational-width campaign independently certified 19 of 64 submitted widths and four registered out-of-sample widths. The latter use $\lambda=5$, parent $p/q=13/21$, and extensions
$(13A+8)/(21A+13)$. Their certified fibre counts are $D$ above the centre, fixing rational IDS $1-D/(21A+13)$.
\begin{table}[ht]
\centering\small
\begin{tabular}{@{}rrl@{}}
\toprule
$A$&$D$&Outward enclosure of $\log(\text{rational gap width})$\\
\midrule
20&2&$[-68.974899318838,-68.974899318837]$\\
20&3&$[-101.720353379131,-101.720353379130]$\\
40&2&$[-67.646792273466,-67.646792273465]$\\
40&3&$[-99.079273854345,-99.079273854344]$\\
\bottomrule
\end{tabular}
\caption{Finite rational certificates at the two registered $D$ values. They are not irrational-transfer certificates.}\label{tab:belowwidthchecks}
\end{table}

For the continuation experiment define the finite-span statistic
\[
 \rho_c=\frac{E_*-\sup\sigma(Q_IH_IQ_I|_{\Ran Q_I})}{m_0}.
\]
The exact dyadic frame coefficients define the tested span.
At $\lambda=5$, parent $13/21$, offset $y=3/10$, the independently certified enclosures are
\[
\begin{array}{ccl}
 b&\text{family}&\rho_c\\
 .28&\text{continued and box}&[66.310252694963,66.323516071840]\\
 .33&\text{continued and box}&[49.271919945713,49.281775315240].
\end{array}
\]
Each family was certified separately; the repeated enclosure does not assert exact equality of their spectra. Every case also passed the strict test $\rho_c>1.000001$. Here the numerical clamp is the specified $C^1$ cubic clamp, whereas the analytic construction uses a fixed $C^\infty$ clamp.

The Green-column campaign independently certified six original rounded-span boxes and two new ones. It measures a cell-boundary tail-rate statistic, rather than the full infinite-volume Schur norm. Its exact window convention and all eight outward intervals are in Numerical Record \S5.

The rational $D=5$ slope registration within $0.25$ of $2(D-1)$ failed in its tested family. The universal corrected-increment residual bound $0.0061$ was also too small; a certified value has magnitude greater than $0.006597244239$. The safe original-sample bound is $0.00660$. In the continuation campaign the original universal no-twofold-drop and finite-$q$ continued/box agreement predictions failed in some preasymptotic cases. The later registered out-of-sample bands passed. These successes do not retroactively change failed registrations and no finite sample proves a limiting constant.

\label{end:diagram}
